\subsection{The image of metric restriction}
\label{subsec:four-point-metric-restriction-image}

The same family allows a precise description of the metrics that do extend. Keep $\mathcal X$, $X_0$, $L$, $\ell$ and the sections $F_1,F_2$ of \cref{sec:four-point-metric-nonextension}. A restricted metric is called genuine if its local weights are not identically $-\infty$. All metric weights below use squared norms and the curvature convention fixed there.

\begin{lemma}\label{lem:psh-harmonic-slice-persistence}
Let $u$ be a genuine plurisubharmonic function on a product of disks $D_\zeta\times D_t$, with $0\in D_t$. Suppose $u(\cdot,t)$ is harmonic for almost every $t$. Then $u(\cdot,0)$ is harmonic or identically $-\infty$.
\end{lemma}
\begin{proof}
Local integrability and Fubini give $\partial_\zeta\partial_{\bar\zeta}u=0$ in distributions. On smaller product disks, form the averages
\[
 u_\varepsilon(\zeta)=\frac1{\pi\varepsilon^2}
       \int_{|t|<\varepsilon}u(\zeta,t)\,\mathrm{d}A(t),
\]
where $\mathrm{d}A$ is planar area. Local upper bounds and reverse Fatou give upper semicontinuity, and Fubini gives the submean inequality in $\zeta$. The averages are locally integrable, since $u$ is integrable on compact product disks. Thus each $u_\varepsilon$ is subharmonic with zero distributional Laplacian, and is harmonic in its canonical representative.

Subharmonicity in $t$ makes these disk averages decrease to $u(\zeta,0)$ as $\varepsilon\downarrow0$. The decreasing limit of harmonic functions is harmonic or identically $-\infty$: subtract the functions from the first one and apply Harnack's inequality and harmonic compactness whenever the values are bounded at one point. This proves the alternative on every smaller disk and hence on $D_\zeta$.
\end{proof}

Let $f:X_0\to\mathbb P^1$ be the resolved pencil of lines through $p_4$. Its affine coordinate is $a=x/y$, and put $P=f^*\mathcal O_{\mathbb P^1}(1)$.

\begin{lemma}\label{lem:four-point-pencil-reference}
The pencil is a holomorphic map with connected fibers, and $P=H-E_4$ and $L_0=\mathcal O(\ell)+P$. Let $h_\ell$ be the divisor metric normalized by $|s_\ell|_{h_\ell}=1$ away from $\ell$, where $s_\ell$ is its canonical section, and let $h_{\mathrm{FS}}$ be the Fubini--Study metric on $\mathcal O_{\mathbb P^1}(1)$, with curvature $\omega_{\mathrm{FS}}$. The metric
\[
 h_*=h_\ell\otimes f^*h_{\mathrm{FS}}
\]
is a genuine restriction of a semipositive metric on $L$ over the whole family. Its curvature is $[\ell]+f^*\omega_{\mathrm{FS}}$.
\end{lemma}
\begin{proof}
First blow up $p_4$ in $\mathbb P^2$. The map $[x:y]$ becomes holomorphic: in the chart $x=r$, $y=rs$, it is $[1:s]$, and the other chart covers the missing direction. Its fibers are the strict lines through $p_4$ and its pencil line is $H-E_4$. Blow up the other three central points. Their directions are $\infty,0,1$, respectively. Each affected fiber becomes the union of its strict transform and one exceptional line meeting it; every other fiber remains a projective line. This proves connectedness and the asserted pullback line bundle. The equality $L_0=\mathcal O(\ell)+P$ follows from the divisor classes.

In this tensor identification, the restrictions of $F_1,F_2$ are $s_\ell f^*x,s_\ell f^*y$. Indeed, at $t=0$ their original forms are $xz,yz$; the pullback of $z=0$ is $\ell+E_1+E_2+E_3$, and removing the four exceptional factors leaves exactly the indicated sections. In a local frame of $L$, the weight $\log(|F_1|^2+|F_2|^2)$ defines a semipositive metric: frame changes multiply both coefficients by the same nonvanishing holomorphic function. The sections have no common zero on any nonzero fiber, so the metric is genuine. Its central weight is the sum of the divisor weight and the pulled-back Fubini--Study weight, proving the claim and the curvature formula.
\end{proof}

\begin{lemma}\label{lem:four-point-variable-conics}
For $a\in\mathbb C\setminus\{0,1\}$, the transformed conics
\[
 F_1-aF_2=z(x-ay)+t(a-1)xy=0
\]
have smooth rational fibers of $L_t$-degree zero when $t\ne0$. Their central fiber is $\ell+q_a$, where $q_a=f^{-1}(a)$ is a smooth projective line of $L_0$-degree one. The two components meet transversely at $p_a=[a:1:0]$, away from the blowup centers. Along $q_a\setminus\{p_a\}$, the conic surface is smooth and its projection to the disk is a submersion.
\end{lemma}
\begin{proof}
The symmetric matrix of the displayed quadratic form has determinant $-a(a-1)t/4$. Its conic is therefore smooth for $t\ne0$ and passes simply through the four centers. Its transformed class is $L_t$, so its degree is $L_t^2=0$. At $t=0$ it splits as the stated two lines; $a\notin\{0,1\}$ and its finiteness exclude the other three central centers from $q_a$. Thus no extra exceptional component occurs, and $q_a=H-E_4$ has degree one against $L_0$. Their intersection is transverse and away from all centers.

Off the centers, the central equation has nonzero spatial derivative along $q_a\setminus\{p_a\}$. At its point on $E_4$, the chart $z=1$, $x=r$, $y=rs$ gives, after removing the exceptional factor,
\[
 1-as+t(a-1)rs=0.
\]
At $r=t=0$, $s=1/a$, its derivative in $s$ is $-a\ne0$. These charts establish the final assertion.
\end{proof}

\begin{theorem}\label{thm:four-point-metric-restriction-image}
The semipositive metrics on $L_0$ that arise as genuine restrictions of semipositive metrics on $L$ over a neighborhood of $X_0$ are exactly
\[
 h_\ell\otimes f^*h_B,
\]
where $h_B$ ranges over all semipositive singular Hermitian metrics on $\mathcal O_{\mathbb P^1}(1)$. Every such central metric extends over the whole family.
\end{theorem}
\begin{proof}
For necessity, shrink the disk by properness so that the extending metric is defined on the whole restricted family, and call its genuine central restriction $h$. On the smooth locus of $f$, local integrability and Fubini show that $h|_{q_a}$ is genuine for almost every $a\in\mathbb C\setminus\{0,1\}$. Fix such an $a$. On each local product chart of the conic surface along $q_a\setminus\{p_a\}$, the extending weight is genuine: its central slice, a restriction of the genuine curve metric, is not identically $-\infty$. Fubini gives genuine metric restrictions on almost every nearby compact conic. Their positive curvature measures have degree zero and hence vanish. \Cref{lem:psh-harmonic-slice-persistence} therefore makes the central weight harmonic on each chart. Its identically $-\infty$ alternative is excluded by the genuine central curve restriction.

The curvature of $h|_{q_a}$ is supported at $p_a$. It is positive of degree one, so it equals $[p_a]$. This is also the curvature of $h_*|_{q_a}$. Their weight difference is a globally defined harmonic function on the projective line: locally at $p_a$, subtracting $\log|w|^2$ from either weight leaves a harmonic function, so their difference extends across that point. It is consequently constant. Thus the weight difference $u$ between $h$ and $h_*$ is constant along almost every fiber, away from $\ell$.

Local holomorphic sections of $f$ avoiding $\ell$ and the exceptional curves exist over every base point. On the finite chart use $[a:1:1]$, and near infinity use $[1:b:1]$ with $b=1/a$. In a local base frame $e_B$, the nonvanishing holomorphic section $s_\ell\otimes f^*e_B$ is a frame of $L_0$ along these sections. Composing the weight of $h$ in this frame with a section defines a subharmonic base weight. It is genuine because the weights on the good fibers are harmonic off their node. Two choices agree for almost every base point because $u$ is constant on those fibers; canonical subharmonic representatives then agree everywhere. Base-frame changes give the metric transition law. These weights therefore define a semipositive metric $h_B$ on $\mathcal O_{\mathbb P^1}(1)$. Fubini on local product charts, with the analytic exceptional fibers and $\ell$ having zero ambient volume, shows that $h$ and $h_\ell\otimes f^*h_B$ agree almost everywhere on $X_0$. Their canonical local psh weights therefore agree everywhere. This proves necessity, including the three exceptional fibers.

For sufficiency, represent $h_B$ by its homogeneous weight $V$ on $\mathbb C^2\setminus\{0\}$. In the chart $Y\ne0$, if $\phi_B$ is the weight in the frame given by the section $Y$, set
\[
 V(X,Y)=\phi_B(X/Y)+\log|Y|^2.
\]
The other chart gives the same function by the metric transition law. It is plurisubharmonic and satisfies $V(\lambda X,\lambda Y)=V(X,Y)+\log|\lambda|^2$ for $\lambda\ne0$. The weight difference between $h_B$ and $h_{\mathrm{FS}}$ is bounded above on the compact base, so
\[
 V(X,Y)\leq\log(|X|^2+|Y|^2)+C.
\]
Thus $V$ extends plurisubharmonically to the origin with value $-\infty$, by \cref{lem:qpsh-removable-analytic}. In each local frame of $L$, the weight $V(F_1,F_2)$ is psh and transforms as a metric weight by homogeneity. It is genuine: the two independent generating sections on each nonzero fiber give a nonconstant map onto $\mathbb P^1$, where $h_B$ is genuine. On the central fiber, the factorization in \cref{lem:four-point-pencil-reference} and homogeneity give exactly the divisor weight plus the pulled-back base weight. This is the required extension over $\mathcal X$.
\end{proof}

\begin{corollary}\label{cor:four-point-least-singular-extendable}
The metric $h_*$ has least singular weights among all genuine extendable restrictions: every such restricted weight $\psi$ satisfies $\psi\leq\psi_*+C$ in compatible local frames, where $\psi_*$ denotes the weight of $h_*$. It is itself extendable. This assertion concerns extendable restrictions, rather than all semipositive metrics on $L_0$.
\end{corollary}
\begin{proof}
By \cref{thm:four-point-metric-restriction-image}, write the metric as $h_\ell\otimes f^*h_B$. Its weight difference from $h_*$ is the pullback of the global quasi-psh difference between $h_B$ and $h_{\mathrm{FS}}$. This function is bounded above on $\mathbb P^1$, proving the inequality. Extendability of $h_*$ is \cref{lem:four-point-pencil-reference}.
\end{proof}
